% !TeX root = ../main.tex
\section{Algebra of DCP operators} \label{appendix:dcp-algebra}

This appendix develops the algebraic tools used to represent and manipulate DCP operators and subgroups.
We first prove the normal form and derive explicit formulas for group operations, then describe the Howell normal form and canonical remainders used to construct canonical subgroup representations.

\subsection{Normal form and group operations}

\begin{dcpnormalformrestatement}[Normal form for DCP operators]
\dcpNormalFormStatement
\end{dcpnormalformrestatement}
\begin{proof}
    We first prove existence by moving all Pauli $X$ factors to
    the right, then prove uniqueness by comparing the action on
    computational-basis states.

    Every DCP operator can be written as a product of scalar
    phases, $\S_i$, $\CZ_{ij}$, and $X_i$, since
    $Z_i=\S_i^2$ and $Y_i=iX_iZ_i$. The relations
    \[
    X_i\S_i=i\S_i^{-1}X_i,
    \qquad
    X_i\CZ_{ij}=Z_j\CZ_{ij}X_i
    \]
    allow us to move each $X_i$ past every diagonal generator;
    generators acting on disjoint qubits commute. The additional
    factors introduced are diagonal, with $Z_j=\S_j^2$, and all
    scalar phases are powers of $\omega_8$. Thus
    \[
    U=\omega_8^k
    \left(\prod_i \S_i^{s_i}\right)
    \left(\prod_{i<j}\CZ_{ij}^{M_{ij}}\right)
    X(\vec{x}).
    \]
    Since $\S_i^4=I$, $\CZ_{ij}^2=I$, and $X_i^2=I$, we may
    take $k\in\Z_8$, $\vec{s}\in\Z_4^n$, and
    $\vec{x}\in\F_2^n$. The exponents $M_{ij}$ uniquely specify
    a symmetric binary matrix with zero diagonal. This proves
    existence.

    To prove uniqueness, suppose
    \[
    \omega_8^k D(\vec{s},M)X(\vec{x})
    =
    \omega_8^{k'}D(\vec{s}',M')X(\vec{x}').
    \]
    Applied to any computational-basis state $\ket{z}$, the two
    sides produce nonzero multiples of $\ket{z\oplus\vec{x}}$
    and $\ket{z\oplus\vec{x}'}$, respectively. Hence
    $\vec{x}=\vec{x}'$, and cancelling the common $X$ factor gives
    \[
    \omega_8^k D(\vec{s},M)
    =
    \omega_8^{k'}D(\vec{s}',M').
    \]
    Both diagonal operators fix $\ket{0^n}$, so evaluating this
    identity on $\ket{0^n}$ gives $k=k'$ in $\Z_8$.
    For the standard basis vector $e_i$,
    \[
    D(\vec{s},M)\ket{e_i}=i^{s_i}\ket{e_i},
    \]
    so $s_i=s_i'$ in $\Z_4$ for every $i$. Finally, for $i<j$,
    \[
    D(\vec{s},M)\ket{e_i\oplus e_j}
    =
    i^{s_i+s_j}(-1)^{M_{ij}}\ket{e_i\oplus e_j}.
    \]
    Since $\vec{s}=\vec{s}'$, this implies $M_{ij}=M_{ij}'$
    for every $i<j$, and therefore $M=M'$. All four parameters
    are unique.
\end{proof}

\begin{proposition}[Group operations on DCP operators]\label{prop:dcp-group-ops}
For $\vec{s}\in\Z_4^{\,n}$ and $M\in\F_2^{\,n\times n}$ symmetric with zero diagonal, write
\begin{align} \label{eq:phi-s-M}
    \varphi_{\vec{s},M}(\vec{v}) \;:=\; \vec{s}\cdot \vec{v} \,+\, 2 \sum_{a<b} M_{ab}\, v_a v_b \pmod 4, \qquad \vec{v} \in \F_2^{\,n},
\end{align}
so that $D(\vec{s},M)\ket{\vec{v}} = i^{\varphi_{\vec{s},M}(\vec{v})}\ket{\vec{v}}$, and let $\odot$ denote the entrywise product of vectors. Given two DCP operators $g_1 = \omega_8^{k_1} D(\vec{s}_1,M_1) X(\vec{x}_1)$ and $g_2 = \omega_8^{k_2} D(\vec{s}_2,M_2) X(\vec{x}_2)$, we have the following formulas for their product, inverse, commutator, and conjugation:
\begin{itemize}
    \item \textbf{Product:}
    \[
    g_1 g_2 \;=\; \omega_8^{\,k_1+k_2+2\varphi_{\vec{s}_2,M_2}(\vec{x}_1)}\, D\big(\vec{s}_1+\vec{s}_2\odot(1-2\vec{x}_1)+2M_2\vec{x}_1,\ M_1+M_2\big)\, X(\vec{x}_1\oplus \vec{x}_2).
    \]
    \item \textbf{Inverse:}
    \[
    g_1^{-1} \;=\; \omega_8^{\,-k_1+2\varphi_{-\vec{s}_1,M_1}(\vec{x}_1)}\, D\big((-\vec{s}_1)\odot(1-2\vec{x}_1)+2M_1\vec{x}_1,\ M_1\big)\, X(\vec{x}_1).
    \]
    \item \textbf{Commutator:}
    \[
    \gcomm{g_1}{g_2} \;:=\; g_1 g_2 g_1^{-1} g_2^{-1} \;=\; i^{\,c}\, Z(\vec{z}),
    \qquad
    c \;=\; \varphi_{\vec{s}_2,M_2}(\vec{x}_1) - \varphi_{\vec{s}_1,M_1}(\vec{x}_2) \pmod 4,
    \]
    \[
    \vec{z} \;=\; M_2\vec{x}_1 + \vec{s}_2\odot \vec{x}_1 + M_1\vec{x}_2 + \vec{s}_1\odot \vec{x}_2 \pmod 2,
    \]
    where $Z(\vec{z}) := Z_1^{z_1} \otimes \dots \otimes Z_n^{z_n}$
    \item \textbf{Conjugation:}
    \[
    g_1 g_2 g_1^{-1} \;=\; \omega_8^{\,k_2+2c}\, D(\vec{s}_2+2\vec{z},\ M_2)\, X(\vec{x}_2),
    \]
    with $(c,\vec{z})$ as in the commutator above.
\end{itemize}
\end{proposition}

\begin{proof}
We first compute how conjugation by $X(\vec{x})$ acts on
a diagonal DCP operator. We then derive the product and
inverse formulas, and obtain the commutator by comparing
the products in the two possible orders.

For binary vectors $\vec{v},\vec{x}$, expanding
$v_a\oplus x_a=v_a+x_a-2v_ax_a$ gives
\begin{align}
    \varphi_{\vec{s},M}(\vec{v}\oplus\vec{x})
    &=
    \varphi_{\vec{s},M}(\vec{x})
    +\varphi_{\vec{t},M}(\vec{v})
    \pmod 4,
    \label{eq:dcp-phase-translation}
    \\
    \vec{t}
    &=
    \vec{s}\odot(1-2\vec{x})+2M\vec{x}
    \pmod 4.
\end{align}
Indeed, the linear terms contribute
$s_a x_a+s_a(1-2x_a)v_a$, while each quadratic term
satisfies
\[
    2M_{ab}(v_a\oplus x_a)(v_b\oplus x_b)
    =
    2M_{ab}
    (v_av_b+v_ax_b+x_av_b+x_ax_b)
    \pmod 4.
\]
Collecting the constant, linear, and quadratic terms gives
\cref{eq:dcp-phase-translation}. Consequently,
\begin{equation}
    X(\vec{x})D(\vec{s},M)X(\vec{x})
    =
    i^{\varphi_{\vec{s},M}(\vec{x})}
    D\big(
        \vec{s}\odot(1-2\vec{x})+2M\vec{x},M
    \big).
    \label{eq:dcp-diagonal-conjugation}
\end{equation}

For the product formula, write
\begin{align*}
    g_1g_2
    &=
    \omega_8^{k_1+k_2}
    D(\vec{s}_1,M_1)
    \big[
        X(\vec{x}_1)D(\vec{s}_2,M_2)X(\vec{x}_1)
    \big]
    X(\vec{x}_1\oplus\vec{x}_2).
\end{align*}
Applying \cref{eq:dcp-diagonal-conjugation} and using
\[
    D(\vec{s},M)D(\vec{s}',M')
    =
    D(\vec{s}+\vec{s}',M+M')
\]
gives the stated expression.

For the inverse formula, note that
$D(\vec{s}_1,M_1)^{-1}=D(-\vec{s}_1,M_1)$. Hence
\[
    g_1^{-1}
    =
    \omega_8^{-k_1}
    \big[
        X(\vec{x}_1)D(-\vec{s}_1,M_1)X(\vec{x}_1)
    \big]
    X(\vec{x}_1),
\]
and another application of
\cref{eq:dcp-diagonal-conjugation} gives the result.

To obtain the commutator, let
\begin{align*}
    \vec{a}_{12}
    &=
    \vec{s}_1+\vec{s}_2\odot(1-2\vec{x}_1)
    +2M_2\vec{x}_1,
    \\
    \vec{a}_{21}
    &=
    \vec{s}_2+\vec{s}_1\odot(1-2\vec{x}_2)
    +2M_1\vec{x}_2
\end{align*}
be the linear parameters in the product formulas for
$g_1g_2$ and $g_2g_1$. Their difference satisfies
\[
    \vec{a}_{12}-\vec{a}_{21}
    =
    2\big(
        M_2\vec{x}_1+\vec{s}_2\odot\vec{x}_1
        +M_1\vec{x}_2+\vec{s}_1\odot\vec{x}_2
    \big)
    =
    2\vec{z}
    \pmod 4.
\]
The two products have the same quadratic parameter
$M_1+M_2$ and the same $X$ parameter
$\vec{x}_1\oplus\vec{x}_2$. Their scalar factors differ
by $i^c$, where
\[
    c=
    \varphi_{\vec{s}_2,M_2}(\vec{x}_1)
    -
    \varphi_{\vec{s}_1,M_1}(\vec{x}_2)
    \pmod 4.
\]
Since $D(2\vec{z},0)=Z(\vec{z})$, it follows that
\[
    g_1g_2=i^c Z(\vec{z})g_2g_1.
\]
Multiplying on the right by $g_1^{-1}g_2^{-1}$ proves
the commutator formula.

Finally,
\begin{align*}
    g_1g_2g_1^{-1}
    &=\gcomm{g_1}{g_2}g_2\\
    &=i^c Z(\vec{z})\,
      \omega_8^{k_2}D(\vec{s}_2,M_2)X(\vec{x}_2)\\
    &=\omega_8^{k_2+2c}
      D(\vec{s}_2+2\vec{z},M_2)X(\vec{x}_2),
\end{align*}
which proves the conjugation formula.
\end{proof}

\subsection{Howell normal form and canonical remainders}
\label{app:howell}

Howell normal form generalizes reduced row-echelon form
from fields to rings $\Z_m$. It provides a canonical
generating matrix for a submodule and a reduction
procedure for choosing representatives of its cosets.
Throughout, we fix the coordinate order and represent
elements of $\Z_m$ by integers in $\{0,\ldots,m-1\}$.

\begin{definition}[Howell normal form]
    Let $N\leq\Z_m^d$. A matrix $H$ with nonzero rows
    $\vec{h}_1,\ldots,\vec{h}_r$ is a \emph{Howell
    normal form} of $N$ if its rows generate $N$ and:
    \begin{enumerate}
        \item The first nonzero entry of $\vec{h}_i$
        occurs in column $j_i$, with
        $j_1<\cdots<j_r$.

        \item Each pivot $p_i:=H_{i,j_i}$ is a positive
        divisor of $m$.

        \item Entries above a pivot satisfy
        $0\leq H_{\ell,j_i}<p_i$ for $\ell<i$.

        \item For every $j\in\{1,\ldots,d\}$,
        \begin{align}
            \{\vec{a}\in N:
              a_1=\cdots=a_{j-1}=0\}
            =
            \operatorname{span}_{\Z_m}
            \{\vec{h}_i:j_i\geq j\}.
            \label{eq:howell-tail}
        \end{align}
    \end{enumerate}
    Zero rows are omitted.
\end{definition}

The last condition ensures that cancellations in early
coordinates are fully captured by the later rows.
For example, over $\Z_8$, the row $(2,1)$ alone is
in echelon form but not Howell form: multiplying it
by $4$ gives $(0,4)$. The Howell normal form of the
module it generates is
\begin{align}
    H=
    \begin{pmatrix}
        2 & 1\\
        0 & 4
    \end{pmatrix}.
\end{align}

Every submodule has a unique Howell normal form.
Given an $s\times d$ generating matrix, this form
can be computed in polynomial time
\cite{storjohann1998fast}.
Thus it depends only on the embedded submodule,
not on its chosen generators. Since pivot positions
are distinct, it has at most $d$ nonzero rows.

The following lemma gives the reduction procedure
used to choose canonical DCP lifts.

\begin{lemma}[Canonical remainders]
    \label{lem:howell-remainders}
    Let $N\leq\Z_m^d$ have Howell normal form $H$,
    with pivot positions $j_1,\ldots,j_r$ and
    pivots $p_1,\ldots,p_r$.
    Every coset of $N$ contains exactly one vector
    $\vec{a}$ satisfying
    \begin{align}
        0\leq a_{j_i}<p_i
        \qquad (i=1,\ldots,r).
        \label{eq:howell-reduced}
    \end{align}
    This vector can be computed by successively
    reducing the pivot coordinates using the rows
    of $H$.
\end{lemma}

\begin{proof}
    We first construct a reduced representative,
    then prove its uniqueness.

    Start with any representative $\vec{a}$ of
    the coset. At step $i$, write its current
    pivot coordinate as
    \begin{align}
        a_{j_i}=q_ip_i+t_i,
        \qquad 0\leq t_i<p_i,
    \end{align}
    and replace $\vec{a}$ by
    $\vec{a}-q_i\vec{h}_i$, with arithmetic performed
    modulo $m$. This preserves the coset and makes
    the $j_i$-th coordinate equal to $t_i$.
    Since $\vec{h}_i$ vanishes before column $j_i$,
    previously reduced pivot coordinates remain
    unchanged. After all rows have been processed,
    \cref{eq:howell-reduced} holds.

    For uniqueness, suppose $\vec{a}$ and $\vec{b}$
    are distinct reduced representatives of the
    same coset. Then
    $\vec{c}:=\vec{a}-\vec{b}\in N$.
    Let $j$ be the first nonzero coordinate of
    $\vec{c}$. By \cref{eq:howell-tail},
    $\vec{c}$ is a linear combination of rows
    whose pivots occur at or after $j$.
    If $j$ were not a pivot position, all such
    rows would vanish in column $j$, a contradiction.
    Hence $j=j_i$ for some $i$, and $c_{j_i}$
    is a multiple of $p_i$ in $\Z_m$.

    Since $p_i$ divides $m$, the integer
    $a_{j_i}-b_{j_i}$ is therefore divisible by
    $p_i$. But both coordinates lie in
    $\{0,\ldots,p_i-1\}$, so
    \begin{align}
        -(p_i-1)
        \leq a_{j_i}-b_{j_i}
        \leq p_i-1.
    \end{align}
    The only multiple of $p_i$ in this interval
    is zero, contradicting the choice of $j$.
    Thus the reduced representative is unique.
\end{proof}

Write $\operatorname{rem}_H(\vec{a})$ for this
unique representative of $\vec{a}+N$.
The lemma immediately gives
\begin{align}
    \vec{a}\in N
    &\iff \operatorname{rem}_H(\vec{a})=\vec{0},
    \\
    \vec{a}+N=\vec{b}+N
    &\iff
    \operatorname{rem}_H(\vec{a})
    =\operatorname{rem}_H(\vec{b}).
\end{align}
Once $H$ is known, reduction takes $O(rd)$
coordinate operations and integer divisions on
representatives smaller than $m$.
In particular, for the fixed modulus $m=8$, it
takes $O(d^2)$ time. The remainder map need not
be linear; it selects representatives of quotient
classes without asserting that the quotient splits.

To apply this construction to a DCP subgroup $G$,
take
\begin{align}
    N=\iota(\Delta(G))\leq\Z_8^d,
    \qquad d=1+n+\binom{n}{2},
\end{align}
and let $H$ be its Howell normal form.
For a lift $h_j=D_jX(\vec{x}^{(j)})\in G$, define
$\widetilde{D}_j$ by
\begin{align}
    \iota(\widetilde{D}_j)
    =
    \operatorname{rem}_H\big(\iota(D_j)\big).
\end{align}
This is a valid diagonal DCP operator: reduction
changes $\iota(D_j)$ by an element of
$\iota(\Delta(G))$, so the result remains in the
image of $\iota$. Moreover,
$\widetilde{D}_jD_j^{-1}\in\Delta(G)$, and hence
\begin{align}
    \widetilde{h}_j
    :=
    \widetilde{D}_jX(\vec{x}^{(j)})
    =
    (\widetilde{D}_jD_j^{-1})h_j
    \in G.
\end{align}
All lifts of $\vec{x}^{(j)}$ have diagonal coordinate
vectors in the same coset of $N$. Therefore
$\widetilde{h}_j$ depends only on $G$, $\vec{x}^{(j)}$,
and the fixed coordinate order, not on the
initial lift.

